\scanpage{527}
\mspage{438}
% B01
\[ x=\varinjlim_I\beta_{i*}(x_i)\qquad \underline{\lim}\ \text{\uuline{dans $E_s$}}. \]

% B02
\deleted{Il faut prouver que} \inserted{On sait que} c'est une $\unsure{\underline{\lim}|\varinjlim}$ dans $E$, donc que pour tt $y$ dans $E$, on a
% B03
\[ \mathrm{Hom}(x,y)\xrightarrow{\sim}\varprojlim_I\mathrm{Hom}(\beta_{i\unsure{0|\circ|*}}(x_i),y) \]
% B04
ou ce qui revient au $=$, pour tte flèche $\lambda:s\unsure{\relbar|\to}t$, et $y$ dans $\mathrm{Ob}\,E_t$, on a
% B05
\[ \begin{aligned}
\mathrm{Hom}^{\unsure{\prime|{}}}_\lambda(x,y)&\ \overset{\sim}{\relbar}\ \varprojlim\mathrm{Hom}_{E_t}(\underbrace{\lambda_*\beta_{i*}(x_i)}_{(\lambda\unsure{{}_i|:|\cdot}\beta_i)_*(x_i)},y)\\
\shortparallel &\hspace{13em}\deleted{\unsure{i}}\rule[.45ex]{6.6em}{.35pt}\rule[.45ex]{.35pt}{.5ex}\\
\mathrm{Hom}_{E_t}(\lambda_*x,y)&
\end{aligned} \]

% B06
Cela résulte \unsure{aussi|alors} du fait que $\lambda_*$ est $\pi$-accessible (étant $\pi_1$-accessible) .... On vérifie facilement que c'est une $\underline{\lim}$ des $x_i$ dans $E$. Et bien sûr, cette li\deleted{\unsure{m}}mite commute avec $f$ par construction.

% B07
Mais dans le cas où $f$ est fibrant, je ne vois comment faire pour construire une limite \unsure{des|dans} \unsure{$E_s$|$E$} -- Il est vrai qu'on peut \inserted{$1^\circ$\unsure{...}} supposer que tous les $s_i$ sont $s$, mais les transitions $\beta_{ji}$ ne sont pas l'identité, ni $=$ des iso. Il faudrait prendre $C=P$ pour essayer de comprendre -- je ne vais pas insister.

% B08
% slanted
\margin{J'ai vérifié \unsure{ici|ce} que si $f$ fibrant, alors $E$ est \uline{\unsure{Kar.}\ \deleted{\unsure{-- donc|donc}}}. Il y a donc déjà \unsure{ça} de \unsure{prouvé|précisé} \unsure{;|.}}

% archive sheet number: 526
